
To determine the expectation value of the fall time, we need to map
  our quantum uncertainties onto classical initial conditions.  We
  model the ice pick as a classical inverted pendulum of length $L$
  and mass $m$.  For small angles ($\sin\theta \approx \theta$), the
  equation of motion is:

  \begin{equation}
    \begin{split}
      I \ddot{\theta} &= m g L \sin\theta \\
      m L^2 \ddot{\theta} &\approx m g L \theta \\
      \ddot{\theta} - \frac{g}{L} \theta &= 0
    \end{split}
  \end{equation}

  This is a standard second-order linear differential equation.
  Defining $\gamma = \sqrt{g/L}$, the general solution is:

  \begin{equation}
    \theta(t) = A e^{\gamma t} + B e^{-\gamma t}
  \end{equation}

  We determine the constants $A$ and $B$ using our initial starting
  angle ($\theta_0$) and initial angular velocity ($\omega_0$):

  \begin{equation}
    \begin{split}
      \theta(0) &= A + B = \theta_0 \\
      \dot{\theta}(0) &= \gamma A - \gamma B = \omega_0
    \end{split}
  \end{equation}

  Solving this system of two equations yields $A = \frac{1}{2}\left(\theta_0 + \frac{\omega_0}{\gamma}\right)$. 
  Because we are looking for a macroscopic fall time, the decaying 
  exponential term ($B e^{-\gamma t}$) vanishes almost instantly, leaving
  us with:

  \begin{equation}
    \theta(t) \approx \frac{1}{2}\left(\theta_0 + \frac{\omega_0}{\gamma}\right) e^{\gamma t}
  \end{equation}

  We define the fall time $t_f$ as the moment the ice pick hits the
  table (e.g., $\theta(t_f) = \pi/2$).  Solving for $t_f$:

  \begin{equation}
    \inlabel{classicalFall}
    \begin{split}
      \frac{\pi}{2} &= \frac{1}{2}\left(\theta_0 + \omega_0\sqrt{\frac{L}{g}}\right) e^{t_f \sqrt{g/L}} \\
      t_f &= \sqrt{\frac{L}{g}} \ln\left( \frac{\pi}{\theta_0 + \omega_0 \sqrt{\frac{L}{g}}} \right)
    \end{split}
  \end{equation}

  Finally, we map our quantum uncertainties directly to these
  classical initial conditions. The position uncertainty $\Delta x$
  becomes our initial angle, and our momentum uncertainty $\Delta p$
  becomes our initial angular velocity:

  \begin{equation}
    \begin{split}
      \theta_0 &\approx \frac{\Delta x}{L} = \frac{d}{L\sqrt{2}} \\
      \omega_0 &\approx \frac{\Delta p}{mL} = \frac{\hbar}{mLd\sqrt{2}}
    \end{split}
  \end{equation}    
